\chapter{Introduction}
\label{chap:introduction}

% STRUCTURE (restructured 2026-09-30, inspired by Nordal (2026), Ch. 1):
%   1.1 Motivation   (absorbs the old "Problem Statement" section)
%   1.2 Background   (absorbs the old Background chapter's literature
%                     sections -- background.tex is no longer \input)
%   1.3 Objective    (absorbs the old "Scope and Delimitations" and
%                     "Report Outline" sections)
% The technical content of the old Background chapter (grasp matrix,
% wrenches, internal forces, LP3) has been moved into the Theory
% chapter; see the provenance comments there.

\section{Motivation}
\label{sec:motivation}

% TODO (from original introduction.tex TODO list, items 1-2):
%  - Why object-aided locomotion matters for snake robots (e.g.
%    traversing environments where pure body-terrain contact is
%    insufficient, obstacle-aided locomotion, etc.).
%  - Problem statement [was its own section "Problem Statement"]:
%    why force closure is the right lens, and why a *metric* (margin)
%    for it is needed for path planning, not just a binary
%    force-closure test.

\section{Background}
\label{sec:background}

% TODO: Situate the project in the literature (cf. Nordal (2026),
% Sec. 1.2 "Previous Work").
There has been done throughout work on form closure for object aided locomotion for snake robots. Utilizing the concept both in path planning and control.

\section{Objective}
\label{sec:objective}

% TODO (from original introduction.tex TODO list, items 3-5):
%  - Objectives / research questions of the project thesis.
%  - Scope and delimitations [was its own section "Scope and
%    Delimitations"]: what is in vs. out of scope (e.g. 2D vs. 3D,
%    specific robot model, specific environment class).
%  - Report outline [was its own section "Report Outline"]: one
%    paragraph summarising each remaining chapter.
